% ===================================================================
% CHAPITRE 1 — ROLE & ARCHITECTURE
% ===================================================================
\chapsec{coverbg}{sitemap}{Chapitre 1 --- R\^ole du Module Administration dans l'\'Ecosyst\`eme MonEcole}

\begin{multicols}{2}

\flyersec{indigo}{bullseye}{Pourquoi ce module~?}

Le module \textbf{Administration} est le \textbf{socle fondateur} de chaque \'etablissement dans MonEcole. Il r\'epond \`a trois questions~:

\begin{enumerate}[label=\textcolor{indigo}{\arabic*.},itemsep=1pt,font=\scriptsize\bfseries]
  \item \textbf{Qui sommes-nous~?} --- Identit\'e, immatriculation, localisation.
  \item \textbf{Comment sommes-nous organis\'es~?} --- Campus, cycles, classes, cours, calendrier.
  \item \textbf{Qui acc\`ede \`a quoi~?} --- Personnel, comptes, droits par module.
\end{enumerate}

\begin{dangerbox}
\scriptsize Sans configuration compl\`ete dans Administration, \textbf{tous les autres modules} (Scolarit\'e, \'Evaluations, Enseignements, Recouvrement) restent \textbf{verrouill\'es}. Le syst\`eme impose un d\'ebloquage progressif.
\end{dangerbox}

\flyersec{hubcolor}{project-diagram}{Architecture Hub-and-Spoke}

MonEcole s\'epare les \textbf{donn\'ees de r\'ef\'erence nationales} des \textbf{donn\'ees op\'erationnelles} de chaque \'ecole~:

\begin{cbox}[title={\faDatabase\ Hub --- \texttt{countryStructure}}]{hubcolor}
\scriptsize R\'ef\'erentiel national~: cycles p\'edagogiques, classes, cours, mentions, grilles de cotation, structure administrative (Provinces, Villes, Communes), calendrier national.
\end{cbox}
\vspace{2pt}
\begin{cbox}[title={\faSchool\ Spoke --- \texttt{db\_monecole}}]{spokecolor}
\scriptsize Donn\'ees op\'erationnelles propres~: \'el\`eves, inscriptions, notes, bulletins, personnel, emplois du temps, frais scolaires.
\end{cbox}

\vspace{2pt}
{\scriptsize\textbf{Le module Administration} est le \textcolor{indigo}{point de jonction} entre Hub et Spoke~: il permet d'activer les cycles/classes du Hub, personnaliser le calendrier, et rattacher l'identit\'e de l'\'ecole \`a la structure nationale.}

\flyersec{violet}{globe}{Acc\`es Multi-Tenant}

Chaque \'etablissement acc\`ede \`a sa propre instance via un \textbf{sous-domaine unique}~:

\begin{center}
\scriptsize\texttt{https://\textcolor{emerald}{\textbf{collegealfajiri}}.monecole.pro}
\end{center}

{\scriptsize Le middleware d\'etecte le sous-domaine, identifie l'\'etablissement dans le Hub, charge son pays/structure/ann\'ee active, et redirige vers le tableau de bord avec les donn\'ees de \textbf{son} \'ecole uniquement. Aucune \'ecole ne peut voir les donn\'ees d'une autre.}

\end{multicols}

% === DIAGRAMME HUB-SPOKE PLEINE LARGEUR ===
\begin{center}
\begin{tikzpicture}[
  hub/.style={draw=hubcolor,fill=hubcolor!6,rounded corners=10pt,minimum width=5cm,minimum height=1.6cm,align=center,font=\scriptsize\bfseries,text=hubcolor,line width=1.2pt},
  spoke/.style={draw=spokecolor,fill=spokecolor!5,rounded corners=7pt,minimum width=2.2cm,minimum height=1cm,align=center,font=\tiny\bfseries,text=spokecolor,line width=.8pt},
  darrow/.style={{Stealth[length=4pt,width=3pt]}-{Stealth[length=4pt,width=3pt]},semithick,color=amber}
]
  \node[hub] (hub) at (0,0) {\faDatabase\hspace{3pt}Hub National\\[-1pt]\tiny\color{slate}countryStructure};
  \node[spoke] (s1) at (-6,-1.8) {\faSchool\ \'Ecole A\\\tiny db\_monecole};
  \node[spoke] (s2) at (-3.2,-1.8) {\faSchool\ \'Ecole B\\\tiny db\_monecole};
  \node[spoke] (s3) at (-0.4,-1.8) {\faSchool\ \'Ecole C\\\tiny db\_monecole};
  \node[spoke] (s4) at (2.4,-1.8) {\faSchool\ \'Ecole D\\\tiny db\_monecole};
  \node[spoke] (s5) at (5.2,-1.8) {\faSchool\ ...N\\\tiny db\_monecole};
  \foreach \s in {s1,s2,s3,s4,s5} \draw[darrow] (hub) -- (\s);
  \node[right,font=\tiny\color{slate}] at (5.5,-0.9) {sync bidirectionnel};
\end{tikzpicture}
\end{center}

% === 6 SECTIONS — OVERVIEW VISUEL ===
\chapsec{slate}{th-large}{Les 6 Sections du Module --- Vue d'Ensemble et D\'ependances}

\begin{multicols}{2}

\flyersec{indigo}{stream}{Flux de d\'ependance}

{\scriptsize Les sections suivent un \textbf{ordre logique}. Chaque section d\'ependant de la pr\'ec\'edente~:}

\begin{center}
\begin{tikzpicture}[
  st/.style={draw=slate!30,fill=white,rounded corners=5pt,minimum width=1.6cm,minimum height=1.1cm,align=center,font=\tiny\bfseries,text=slate,drop shadow={shadow xshift=.5pt,shadow yshift=-.5pt,opacity=0.08}},
  ar/.style={-{Stealth[length=4pt,width=3pt]},semithick,color=slatelight}
]
  \node[st,fill=indigo!6,draw=indigo!30] (a) at (0,0) {\faThLarge\\Dashboard};
  \node[st,fill=sky!6,draw=sky!30,right=2pt of a] (b) {\faSchool\\Mon \'Ecole};
  \node[st,fill=emerald!6,draw=emerald!30,right=2pt of b] (c) {\faCogs\\Config};
  \node[st,fill=amber!6,draw=amber!30,below=4pt of a] (d) {\faCalendar\\Calendrier};
  \node[st,fill=amber!6,draw=amber!30,below=4pt of b] (e) {\faUsers\\Personnel};
  \node[st,fill=rose!6,draw=rose!30,below=4pt of c] (f) {\faIcon{shield-alt}\\Utilisateurs};
  \draw[ar] (a)--(b); \draw[ar] (b)--(c);
  \draw[ar] (c)--(d); \draw[ar] (c)--(e); \draw[ar] (e)--(f);
\end{tikzpicture}
\end{center}

\flyersec{rose}{lock}{Verrouillage progressif}

{\scriptsize Le syst\`eme \textbf{verrouille} les sections avanc\'ees tant que les pr\'erequis ne sont pas remplis~:}

{\scriptsize
\begin{tabularx}{\columnwidth}{lX}
\toprule
\textbf{Section verrouill\'ee} & \textbf{Condition de d\'eblocage} \\
\midrule
Configuration & G\'eolocalisation compl\`ete \\
Calendrier & G\'eolocalisation compl\`ete \\
Personnel & G\'eolocalisation compl\`ete \\
Utilisateurs & G\'eolocalisation compl\`ete \\
\textit{Onglet Cycles} & $\geq$1 campus cr\'e\'e \\
\textit{Onglet Classes} & $\geq$1 campus cr\'e\'e \\
\bottomrule
\end{tabularx}
}

\end{multicols}

% === IMPACT SUR LES AUTRES MODULES ===
\begin{center}
\begin{tikzpicture}[
  mod/.style={draw=#1!50,fill=#1!6,rounded corners=5pt,minimum width=3cm,minimum height=.7cm,font=\tiny\bfseries,text=slate,align=center},
  dep/.style={-{Stealth[length=3pt]},semithick,color=#1!50}
]
  \node[mod=indigo] (admin) at (0,0) {\faCogs\ Administration};
  \node[mod=emerald] (scol) at (-5.5,-1.3) {\faUserGraduate\ Scolarit\'e};
  \node[mod=violet] (ens) at (-1.8,-1.3) {\faChalkboardTeacher\ Enseignements};
  \node[mod=sky] (eval) at (1.8,-1.3) {\faClipboard\ \'Evaluations};
  \node[mod=amber] (rec) at (5.5,-1.3) {\faMoneyBill\ Recouvrement};
  \draw[dep=emerald] (admin)--(scol) node[midway,above,font=\tiny\color{slate},sloped]{classes, campus};
  \draw[dep=violet] (admin)--(ens) node[midway,above,font=\tiny\color{slate},sloped]{personnel};
  \draw[dep=sky] (admin)--(eval) node[midway,above,font=\tiny\color{slate},sloped]{p\'eriodes, cours};
  \draw[dep=amber] (admin)--(rec) node[midway,above,font=\tiny\color{slate},sloped]{classes, \'el\`eves};
  \draw[dep=violet,dashed] (ens)--(eval) node[midway,below,font=\tiny\color{slate}]{attribution};
\end{tikzpicture}
\end{center}

\newpage
